% O014 / D80 -- Methods of Algebra, Volume 2 -- independent English edition
% Source: chapter9.tex, lines 1724--1804 inclusive.
% stable-unit-id: o014.aljabr2.chapter9.fiber-functor-isomorphisms
% Verified source corrections carried from the Indonesian critical witness:
% omega'_1 on source line 1753 is omega_1; omega_2(X) on source line 1767 is
% omega_1(Y), as required by the surrounding diagram.
Finally, we briefly discuss the relationship between fiber functors. Unlike
the preceding discussion, this part does not involve the reconstruction
theorem. Let $\mathcal{T}$ be a Tannakian category, and let
$\omega_1,\omega_2$ be fiber functors over a commutative $\Bbbk$-algebra
$B$. For every commutative $B$-algebra $B'$, define the set
\begin{equation*}
	\iHom^\otimes_B(\omega_2, \omega_1)(B') := \Hom_{\text{monoidal functors}}\left( \omega_2 \dotimes{B} B', \omega_1 \dotimes{B} B' \right).
\end{equation*}
As $B'$ varies, this defines a functor
$B\dcate{CAlg}\to\cate{Set}$. Recall that
Proposition~\ref{prop:L-bialgebra}~(iii) equips
$\coHom_B(\omega_1,\omega_2)$ with the structure of a commutative
$B$-algebra.

\begin{proposition}\label{prop:Isom-Tannakian}
	For a Tannakian category $\mathcal{T}$ and fiber functors
	$\omega_1,\omega_2:\mathcal{T}\to\cated{Mod}B$, there is an isomorphism
	of functors
	\[ \iHom^\otimes_B(\omega_2, \omega_1) \simeq \Hom_{B\dcate{CAlg}}\left( \coHom_B(\omega_1, \omega_2) , \; \cdot \; \right). \]
	Composition of homomorphisms on the right is induced by
	\eqref{eqn:L-comult-prep}.
\end{proposition}
\begin{proof}
	To simplify notation, write $\coHom_B(\omega_1,\omega_2)$ as $\coHom$
	from now on. Specifying a morphism, without regard to the monoidal
	structure,
	\[ \tilde{\varphi}: \omega_2 \dotimes{B} B' \to \omega_1 \dotimes{B} B' \quad (\text{both valued in }\cated{Mod}B') \]
	is equivalent to specifying a morphism
	\[ \varphi: \omega_2 \to \omega_1 \dotimes{B} B' \quad (\text{both valued in }\cated{Mod}B), \]
	which, by the universal property of $\coHom$, is in turn equivalent to
	specifying a homomorphism of $B$-modules
	\[ \psi: \coHom \to B'. \]
	It remains to prove that $\tilde{\varphi}$ preserves the monoidal
	structure if and only if $\psi$ is a homomorphism of $B$-algebras. The
	argument is formal, although spelling it out is somewhat lengthy.

	Write the multiplication morphisms as
	$m:B'\dotimes{B}B'\to B'$ and
	$\mu:\coHom\dotimes{B}\coHom\to\coHom$. It is immediate that
	$\tilde{\varphi}$ preserves the unit, respectively $\otimes$, if and only
	if the left-hand, respectively right-hand, diagram below always commutes:
	\[\begin{tikzcd}
		\omega_2(\munit) \arrow[dd, "\sim"' sloped] \arrow[r, "{\varphi_{\munit}}"] & \omega_1(\munit) \dotimes{B} B' \arrow[d, "\sim" sloped] \\
		& B \dotimes{B} B' \arrow[d, "\sim" sloped] \\
		B \arrow[r] & B'
	\end{tikzcd} \quad
	\begin{tikzcd}
		\omega_2(X \otimes Y) \arrow[r, "{\varphi_{X \otimes Y}}"] \arrow[dd, "\sim"' sloped] & \omega_1(X \otimes Y) \dotimes{B} B' \arrow[d, "\sim" sloped] \\
		& \omega_1(X) \dotimes{B} \omega_1(Y) \dotimes{B} B' \\
		\omega_2(X) \dotimes{B} \omega_2(Y) \arrow[r, "{\varphi_X \otimes \varphi_Y}"'] & \omega_1(X) \dotimes{B} B' \dotimes{B} \omega_1(Y) \dotimes{B} B' . \arrow[u, "{\text{induced by }m}"']
	\end{tikzcd}\]

	First consider the left-hand diagram. Its first row factors as
	$\omega_2(\munit)\to\omega_1(\munit)\dotimes{B}\coHom
	\xrightarrow{\identity\otimes\psi}\omega_1(\munit)\dotimes{B}B'$.
	After identifying every $\omega_i(\munit)$ with $B$, commutativity of the
	left-hand diagram is equivalent to commutativity of
	$\begin{tikzcd}[column sep=small, row sep=small]
		B \arrow[r] \arrow[d] & B' \\
		\coHom \arrow[ru, "\psi"'] &
	\end{tikzcd}$,
	which in turn is equivalent to $\psi$ preserving the unit.

	We now show that commutativity of the right-hand diagram is equivalent to
	$\psi$ preserving multiplication. Consider
	\begin{equation}\label{eqn:Isom-Tannakian-aux0}
		\begin{tikzcd}
			\omega_2(X \otimes Y) \arrow[rr, bend left, "{\varphi_{X \otimes Y}}"] \arrow[r] \arrow[dd, "\sim"' sloped] & \omega_1(X \otimes Y) \dotimes{B} \coHom \arrow[r, "{\identity \otimes \psi}"] & \omega_1(X \otimes Y) \dotimes{B} B' \arrow[d, "\sim" sloped] \\
			& \omega_1(X \otimes Y) \dotimes{B} \coHom \dotimes{B} \coHom \arrow[u, "{\identity \otimes \mu}"] & \omega_1(X) \dotimes{B} \omega_1(Y) \dotimes{B} B' \\
			\omega_2(X) \dotimes{B} \omega_2(Y) \arrow[rr, bend right, "{\varphi_X \otimes \varphi_Y}"'] \arrow[r] & \omega_1(X) \dotimes{B} \coHom \dotimes{B} \omega_1(Y) \dotimes{B} \coHom \arrow[r, "{(\identity \otimes \psi)^{\otimes 2}}" inner sep=0.6em] \arrow[u, "\sim" sloped] & \omega_1(X) \dotimes{B} B' \dotimes{B} \omega_1(Y) \dotimes{B} B' , \arrow[u, "{\text{induced by }m}"']
		\end{tikzcd}
	\end{equation}
	The upper and lower curved regions always commute, while commutativity of
	the square on the left reduces to the characterization of $\mu$ in
	\eqref{eqn:L-bialgebra-prod}.

	Define $\eta:=\omega_1^\vee\dotimes{B}\omega_2$, which is a monoidal
	functor $\mathcal{T}\to\cated{Mod}B$. The square portion of
	\eqref{eqn:Isom-Tannakian-aux0} can be rearranged as
	\begin{equation}\label{eqn:Isom-Tannakian-aux}
		\begin{tikzcd}
			\eta(X \otimes Y) \arrow[r] \arrow[d, "\sim"' sloped] & \coHom \arrow[r, "\psi"] & B' \\
			\eta(X) \dotimes{B} \eta(Y) \arrow[r] & \coHom \dotimes{B} \coHom \arrow[r, "{\psi \otimes \psi}"'] \arrow[u, "\mu"] & B' \dotimes{B} B' , \arrow[u, "m"']
		\end{tikzcd}
	\end{equation}
	The left-hand square still commutes, while $\psi$ preserves multiplication
	if and only if the right-hand square commutes.

	If $\tilde{\varphi}$ preserves $\otimes$, then the large outer frame in
	\eqref{eqn:Isom-Tannakian-aux0} commutes, and hence so does the outer frame
	of its square portion. Thus the outer frame of
	\eqref{eqn:Isom-Tannakian-aux} commutes. This shows that the right-hand
	square commutes after precomposition with
	$\eta(X)\dotimes{B}\eta(Y)\to\coHom\dotimes{B}\coHom$. Since $X$ and
	$Y$ are arbitrary, the universal property of $\coHom$, or its concrete
	construction in Proposition~\ref{prop:cogebra-L-universal}, suffices to
	show that the right-hand square in \eqref{eqn:Isom-Tannakian-aux}
	commutes.

	Conversely, suppose the right-hand square in
	\eqref{eqn:Isom-Tannakian-aux} commutes. Then the entire diagram
	\eqref{eqn:Isom-Tannakian-aux}
	commutes, so the large outer frame in \eqref{eqn:Isom-Tannakian-aux0}
	commutes; hence $\tilde{\varphi}$ preserves $\otimes$. This completes the
	verification.
\end{proof}

The rigidity of $\mathcal{T}$ implies that every element of
$\iHom^\otimes_B(\omega_2,\omega_1)(B')$ is an isomorphism; this is an
application of Proposition~\ref{prop:monoidal-functor-rigidity}. The next
result is therefore immediate.

\begin{corollary}\label{prop:Aut-Tannakian}
	\index[sym1]{Aut-tensor@$\Aut^{\otimes}$}
	For a Tannakian category $\mathcal{T}$, a commutative $\Bbbk$-algebra
	$B$, and a fiber functor $\omega:\mathcal{T}\to\cated{Mod}B$, write
	$\Aut^{\otimes}(\omega)$ for the group of monoidal automorphisms of
	$\omega$. Then, for every commutative $B$-algebra $B'$, there is a
	canonical isomorphism
	\[ \Aut^{\otimes}\left( \omega \dotimes{B} B' \right) \simeq \Hom_{B\dcate{CAlg}}\left(\coEnd_B(\omega), B' \right), \]
	where multiplication in the group on the right is reflected by the
	comultiplication of $\coEnd_B(\omega)$.
\end{corollary}

One can further show that the identity element and inversion in the
automorphism group arise respectively from the counit morphism and the
antipode of $\coEnd_B(\omega)$. The verification is simple and formal, so it
is left as an exercise for this chapter. Considering the Yoneda embedding of
$B\dcate{CAlg}$ gives the following conclusion:
\[\begin{tikzcd}
	\text{knowing the functor }\Aut^{\otimes}\left( \omega \dotimes{B} (\cdot)\right): B\dcate{CAlg} \to \cate{Grp} \arrow[Leftrightarrow, d] \\
	\text{knowing the Hopf $B$-algebra }\coEnd_B(\omega).
\end{tikzcd}\]
Thus, once $\Aut^{\otimes}$ is viewed functorially, the entire long story of
the endomorphism coalgebra ultimately returns to the group of
$\otimes$-automorphisms of the fiber functor.

In the situation of Proposition~\ref{prop:Isom-Tannakian}, Deligne proved in
\cite[1.11--1.13]{Del90} that $\coHom_B(\omega_1,\omega_2)$ is faithfully
flat as a $B$-module; in particular,
$\coHom_B(\omega_1,\omega_2)\neq0$. We shall not give the detailed proof,
but its importance in algebraic geometry is worth mentioning. It guarantees
the existence of a faithfully flat $B$-algebra $B'$ and an isomorphism
$\omega_1\dotimes{B}B'\simeq\omega_2\dotimes{B}B'$; for example, one may
take $B'=\coHom_B(\omega_1,\omega_2)$. Thus fiber functors are unique up to
faithfully flat change of rings, while properties over the original ring $B$
can in principle be reconstructed from descent data as in
\S\ref{sec:descent-modules}. This is the principal benefit of faithful
flatness.
